\scanpage{157}
\mspage{73}
% archive sheet number: 156

mais c'est

\[
\begin{aligned}
\mathrm{Filt}^{\prime\,\pi}M
&=\text{ss-catégorie pleine de }M\text{ formée}\\
&\phantom{{}={}}\text{des objets }x\text{ tels qu'il existe}\\
&\phantom{{}={}}\text{une famille épimorphique}\\
&\phantom{{}={}}\text{stricte }\ x_i\xrightarrow{\alpha_i}x\ \text{de cardinal}\\
&\phantom{{}={}}\leq\pi,\ \text{avec }x_i\in S\quad\forall\,i\in I.
\end{aligned}
\tag*{(4.3.6.3)}
\]

\unsure{Aussitôt|Aisément|A priori} on a

\[
\mathrm{Filt}^\pi M\subset\mathrm{Filt}^{\prime\,\pi}M,
\]

mais dans loc. cit on prouve que pour $x$ dans $\mathrm{Filt}^{\prime\,\pi}M$, on a

\[
\mathrm{card}\,\mathrm{Fl}(S_{/x})\leq\pi^{\pi_0},
\]

d'où

\[
\mathrm{Filt}^\pi M\subset\mathrm{Filt}^{\prime\,\pi}M\subset\mathrm{Filt}^{\pi^{\pi_0}}M
\]

donc si $\pi=\pi^{\pi_0}$, p. ex. $\pi=2^c$ avec $c\geq\pi_0$, donc on trouve

\[
\mathrm{Filt}^\pi M=\mathrm{Filt}^{\prime\,\pi}M.
\]

Donc les deux filtrations ont l'air $\pm$ équivalentes. \deleted{Pour n'avoir pas à refaire du travail de SGA 4 I 9.13, je vais essayer de travailler avec la filtration définie dans loc. cit., plutôt qu'avec celle indiquée d'abord.}

Mais la démonstration donnée de la \deleted{condition b) de SGA 4 I 9.12.} stabilité de $\mathrm{Filt}^{\prime\,\pi}$ \deleted{fait} semble vraiment faire appel de façon essentielle à la condition d'universalité des familles épimorphiques strictes, (\deleted{\unsure{donc|dans}} \inserted{\unsure{opéra-|à prio-}} \unsure{attention|portion}, \unsure{\}} car les hypothèses faites \deleted{\unsure{déjà}} impliquent l'existence de petites $\underleftarrow{\lim}$).

Donc \unsure{il|zut}, me faut \unsure{ben} oui, \unsure{mal gré|malgré|mal que} refaire le travail !

